\documentclass[12pt]{article}
\usepackage[utf8]{inputenc}
\usepackage[T1]{fontenc}
\usepackage{hyperref}
\title{Siang Language Resources for K-12 Curriculum Development}
\author{Compiled by Claude Sonnet 5.5 (claude-sonnet-5-5) for Andr\'as Kornai}
\date{}
\begin{document}
\maketitle

\section{Introduction}
Siang (Ot Siang; ISO 639-3 sya) is an Austronesian language of Central Kalimantan, Indonesia, in Murung Raya regency. It is closely related to Murung. Sources disagree on its classification and speaker numbers: Wikipedia gives 60,000 (1981), Joshua Project 20,000. This report rests on a single web search, and little material was found.

\section{Primary Educational Resources}
No primer, textbook or school programme was found in the search results.

\section{Linguistic and Reference Materials}
\begin{itemize}
\item Siang language, Wikipedia (includes a Siang--Murung vocabulary comparison): \url{https://en.wikipedia.org/wiki/Siang_language}
\item Languages of Kalimantan, Wikipedia: \url{https://en.wikipedia.org/wiki/Languages_of_Kalimantan}
\item Endangered Languages Project entry: \url{https://endangeredlanguages.com/lang/5412?hl=pt-br}
\end{itemize}
No dictionary or grammar was found.

\section{Religious and Community Resources}
Joshua Project, Siang in Indonesia, reports that a New Testament and the Jesus Film exist in Siang: \url{https://joshuaproject.net/people_groups/14933}

\section{Digital and Audio Resources}
Global Recordings Network lists recordings in Siang: \url{https://globalrecordings.net/en/language/sya}

\section{Curriculum Development Context}
Search results gave no evidence of Siang being taught in schools. Indonesian is the school language. A curriculum would need new materials and collaboration with Siang-Murung communities. Readers should confirm the New Testament's availability and source separately.
\end{document}